\documentclass[11pt]{article}

% ============================================================
% Packages
% ============================================================
\usepackage[margin=1in]{geometry}
\usepackage{amsmath,amssymb}
\usepackage{array}
\usepackage{booktabs}
\usepackage{enumitem}
\usepackage[T1]{fontenc}
\usepackage{lmodern}
\usepackage{microtype}
\usepackage{graphicx}

% ============================================================
% General Formatting
% ============================================================
\setlength{\parindent}{0pt}
\setlength{\parskip}{0.65em}

\setlist[itemize]{leftmargin=1.5em}
\setlist[enumerate]{leftmargin=1.75em}

\renewcommand{\arraystretch}{1.15}

% Screenshots are assumed to be in the same directory as this .tex file.
\graphicspath{{./}}

% ============================================================
% WHITEBOARD Environment
% ============================================================
\newenvironment{whiteboard}
{
    \par\medskip
    \hrule
    \vspace{0.35em}
    \noindent\textbf{\large WHITEBOARD}
    \par\vspace{0.3em}
    \hrule
    \vspace{0.65em}
}
{
    \vspace{0.45em}
    \hrule
    \par\medskip
}

% ============================================================
% Textbook Reference Helper
% ============================================================
\newcommand{\corresponds}[1]{%
    \textit{Corresponds to Mano, Fig.~#1.}\par\smallskip
}

% ============================================================
% Document
% ============================================================
\begin{document}

\begin{center}
    {\LARGE\textbf{CS2614 --- Computer System Architecture}}\\[0.45em]
    {\Large\textbf{Chapter 2: Digital Components}}\\[0.45em]
    \textit{Based on M. Morris Mano,
    Computer System Architecture, 3rd ed.}
\end{center}

% ============================================================
\section{Integrated Circuits}
% ============================================================

Digital circuits are constructed with integrated circuits.

An integrated circuit, or IC, is a small semiconductor chip containing
electronic components interconnected to perform one or more digital
functions.

As integrated-circuit technology developed, progressively larger numbers of
gates could be placed within a single package.

\begin{whiteboard}

\[
\boxed{
\text{IC}
=
\text{interconnected digital circuitry on one chip}
}
\]

\[
\begin{array}{c|l}
\text{SSI}
&
\text{several gates}
\\
\text{MSI}
&
\text{approximately }10\text{--}200\text{ gates}
\\
\text{LSI}
&
\text{hundreds to thousands of gates}
\\
\text{VLSI}
&
\text{thousands of gates and beyond}
\end{array}
\]

\end{whiteboard}

Integrated circuits are also classified according to their digital logic
family.

Mano identifies TTL, ECL, MOS, and CMOS as major families.

TTL became a widely used standard family.

ECL is associated with high-speed operation.

MOS technology permits high component density.

CMOS is particularly useful when low power consumption is important.

\begin{whiteboard}

\[
\boxed{\text{Digital Logic Families}}
\]

\[
\begin{array}{c|l}
\text{TTL}
&
\text{widely used standard family}
\\
\text{ECL}
&
\text{high speed}
\\
\text{MOS}
&
\text{high component density}
\\
\text{CMOS}
&
\text{low power consumption}
\end{array}
\]

\end{whiteboard}

The remainder of the chapter introduces digital components that serve as
reusable building blocks in computer organization.

% ============================================================
\section{Decoders}
% ============================================================

A binary code of \(n\) bits can represent as many as

\[
2^n
\]

distinct combinations.

A decoder is a combinational circuit that converts an \(n\)-bit coded input
into as many as \(2^n\) distinct output lines.

A decoder with \(n\) inputs and \(m\) outputs is commonly called an
\(n\)-to-\(m\)-line decoder, where

\[
m\leq 2^n.
\]

\begin{whiteboard}

\[
\boxed{
n\text{-to-}m\text{ decoder},
\qquad
m\leq2^n
}
\]

For the full decoding case:

\[
\boxed{
n\text{ inputs}
\rightarrow
2^n\text{ outputs}
}
\]

\end{whiteboard}

% ------------------------------------------------------------
\subsection{3-to-8-Line Decoder}
% ------------------------------------------------------------

A 3-to-8-line decoder accepts three binary inputs and produces eight output
lines.

Each output corresponds to one of the eight possible combinations of the
three input variables.

When enabled, exactly one output is active for a given input combination.

\begin{whiteboard}

\corresponds{2-1}

\begin{center}
\includegraphics[
    width=0.90\textwidth
]{1_fig_2_1.png}
\end{center}

\[
A_2A_1A_0
\rightarrow
D_0,D_1,\ldots,D_7
\]

\[
\boxed{
\text{one output corresponds to each 3-bit combination}
}
\]

\end{whiteboard}

The decoder can also be interpreted as a binary-to-octal converter.

For example,

\[
A_2A_1A_0=110
\]

selects output

\[
D_6.
\]

\begin{whiteboard}

\[
A_2A_1A_0=110_2=6_{10}
\]

\[
\boxed{D_6=1}
\]

All other decoder outputs:

\[
D_i=0
\qquad
(i\neq6)
\]

\end{whiteboard}

% ------------------------------------------------------------
\subsection{Enable Input}
% ------------------------------------------------------------

Commercial decoders normally contain one or more enable inputs.

For the decoder of Fig.~2-1, the circuit is enabled when

\[
E=1.
\]

When

\[
E=0,
\]

all outputs are 0 regardless of the values of the three data inputs.

\begin{whiteboard}

\[
\boxed{
E=0
\Rightarrow
D_0=D_1=\cdots=D_7=0
}
\]

\[
\boxed{
E=1
\Rightarrow
\text{normal decoding}
}
\]

When enabled:

\[
\boxed{
\text{exactly one output}=1
}
\]

\end{whiteboard}

% ------------------------------------------------------------
\subsection{NAND-Gate Decoder and Active-Low Logic}
% ------------------------------------------------------------

A decoder may also be constructed using NAND gates.

In this arrangement the outputs may be generated in complemented form.

The decoder of Fig.~2-2 is enabled when its complemented enable input is 0.
When enabled, exactly one output is 0 and the remaining outputs are 1.

\begin{whiteboard}

\corresponds{2-2}

\begin{center}
\includegraphics[
    width=0.90\textwidth
]{2_fig_2_2.png}
\end{center}

\[
\boxed{
\text{active-high signal: asserted}=1
}
\]

\[
\boxed{
\text{active-low signal: asserted}=0
}
\]

For this decoder:

\[
\boxed{
\bar E=0
\Rightarrow
\text{enabled}
}
\]

\[
\boxed{
\text{exactly one output}=0
}
\]

\end{whiteboard}

% ------------------------------------------------------------
\subsection{Decoder Expansion}
% ------------------------------------------------------------

Enable inputs make it possible to combine smaller decoders into a larger
decoder.

Mano demonstrates this by constructing a 3-to-8-line decoder from two
2-to-4-line decoders.

The two least significant input bits are applied to both smaller decoders.

The most significant bit determines which decoder is enabled.

\begin{whiteboard}

\corresponds{2-3}

\begin{center}
\includegraphics[
    width=0.66\textwidth
]{3_fig_2_3.png}
\end{center}

\[
A_1,A_0
\rightarrow
\text{both 2-to-4 decoders}
\]

\[
A_2=0
\Rightarrow
D_0\text{--}D_3
\]

\[
A_2=1
\Rightarrow
D_4\text{--}D_7
\]

\[
\boxed{
\text{Enable inputs allow component expansion}
}
\]

\end{whiteboard}

% ------------------------------------------------------------
\subsection{Encoders}
% ------------------------------------------------------------

An encoder performs the inverse operation of a decoder.

A decoder converts a binary code into one of many outputs.

An encoder converts one active input line into a binary code.

\begin{whiteboard}

\[
\boxed{
\text{Decoder: code}
\rightarrow
\text{one-of-many}
}
\]

\[
\boxed{
\text{Encoder: one-of-many}
\rightarrow
\text{code}
}
\]

\end{whiteboard}

For an octal-to-binary encoder, eight input lines correspond to the octal
digits 0 through 7 and three output lines produce the equivalent binary code.

The circuit assumes that only one input is active at a time.

The output equations are

\[
A_0=D_1+D_3+D_5+D_7,
\]

\[
A_1=D_2+D_3+D_6+D_7,
\]

\[
A_2=D_4+D_5+D_6+D_7.
\]

\begin{whiteboard}

\[
\boxed{
8\text{-to-}3\text{ encoder}
}
\]

\[
\boxed{
\text{Only one input is assumed to be }1
}
\]

\[
A_0=D_1+D_3+D_5+D_7
\]

\[
A_1=D_2+D_3+D_6+D_7
\]

\[
A_2=D_4+D_5+D_6+D_7
\]

\end{whiteboard}

% ============================================================
\section{Multiplexers}
% ============================================================

A multiplexer is a combinational circuit that selects binary information from
one of several data input lines and directs the selected value to a single
output.

A \(2^n\)-to-1 multiplexer contains

\[
2^n
\]

data input lines and

\[
n
\]

selection inputs.

The multiplexer is therefore also called a \emph{data selector}.

\begin{whiteboard}

\[
\boxed{
2^n\text{-to-}1\text{ MUX}
}
\]

\[
2^n\text{ data inputs}
\]

\[
n\text{ selection inputs}
\]

\[
1\text{ output}
\]

\[
\boxed{\text{MUX}=\text{data selector}}
\]

\end{whiteboard}

% ------------------------------------------------------------
\subsection{4-to-1 Multiplexer}
% ------------------------------------------------------------

A 4-to-1 multiplexer has four data inputs,

\[
I_0,I_1,I_2,I_3,
\]

two selection inputs,

\[
S_1,S_0,
\]

and one output,

\[
Y.
\]

\begin{whiteboard}

\corresponds{2-4}

\begin{center}
\includegraphics[
    width=0.82\textwidth
]{4_fig_2_4.png}
\end{center}

\[
\begin{array}{cc|c}
S_1&S_0&Y\\
\hline
0&0&I_0\\
0&1&I_1\\
1&0&I_2\\
1&1&I_3
\end{array}
\]

\end{whiteboard}

For example,

\[
S_1S_0=10
\]

selects input

\[
I_2.
\]

Thus,

\[
Y=I_2.
\]

\begin{whiteboard}

\[
S_1S_0=10
\]

\[
\boxed{
Y=I_2
}
\]

\[
\boxed{
\text{selection lines determine which input reaches }Y
}
\]

\end{whiteboard}

% ============================================================
\section{Registers}
% ============================================================

A register is a group of flip-flops used to store binary information.

Since each flip-flop stores one bit, an \(n\)-bit register contains \(n\)
flip-flops and can store one \(n\)-bit binary quantity.

\begin{whiteboard}

\[
\boxed{
n\text{-bit register}
=
n\text{ flip-flops}
}
\]

\[
\boxed{
1\text{ flip-flop}
=
1\text{ stored bit}
}
\]

\[
\boxed{
\text{register stores one binary word}
}
\]

\end{whiteboard}

% ------------------------------------------------------------
\subsection{4-Bit Register}
% ------------------------------------------------------------

The simplest register consists of flip-flops sharing a common clock.

In Mano's 4-bit register, four D flip-flops store four input bits.

At the active clock transition, all four inputs are transferred into the
register simultaneously.

This operation is called a \emph{parallel load}.

\begin{whiteboard}

\corresponds{2-6}

\begin{center}
\includegraphics[
    height=0.58\textheight,
    keepaspectratio
]{5_fig_2_6.png}
\end{center}

\[
I_0,I_1,I_2,I_3
\rightarrow
A_0,A_1,A_2,A_3
\]

\[
\boxed{
\text{all bits load simultaneously}
}
\]

\[
\boxed{
\text{simultaneous loading}
=
\text{parallel load}
}
\]

\end{whiteboard}

The clear input resets all flip-flops asynchronously.

During normal clocked operation, the clear input remains inactive.

% ------------------------------------------------------------
\subsection{Register with Parallel Load Control}
% ------------------------------------------------------------

A computer normally has a continuously running master clock.

A separate control signal is therefore needed to determine whether a
particular register should accept new information.

In the register of Fig.~2-7, the clock is continuously available while the
Load input determines what appears at each D input.

When

\[
\text{Load}=1,
\]

the external input is selected.

When

\[
\text{Load}=0,
\]

the current register output is fed back to the D input so that the state is
preserved.

\begin{whiteboard}

\corresponds{2-7}

\begin{center}
\includegraphics[
    height=0.62\textheight,
    keepaspectratio
]{6_fig_2_7.png}
\end{center}

\[
\boxed{
\text{Load}=1
\Rightarrow
D_i=I_i
}
\]

\[
\boxed{
\text{Load}=0
\Rightarrow
D_i=A_i
}
\]

\[
\boxed{
\text{feedback provides the hold condition}
}
\]

\end{whiteboard}

The feedback path is necessary because the D flip-flop itself has no separate
``no change'' input combination.

% ============================================================
\section{Shift Registers}
% ============================================================

A shift register is a register capable of shifting its binary information from
one stage to the next.

The output of each flip-flop is connected to the D input of the following
flip-flop.

All stages share a common clock.

\begin{whiteboard}

\[
\boxed{
\text{Shift register}
=
\text{register that moves stored bits}
}
\]

\[
\text{Serial In}
\rightarrow
A_0
\rightarrow
A_1
\rightarrow
A_2
\rightarrow
A_3
\rightarrow
\text{Serial Out}
\]

\[
\boxed{
1\text{ clock transition}
\Rightarrow
1\text{ shift}
}
\]

\end{whiteboard}

% ------------------------------------------------------------
\subsection{4-Bit Shift Register}
% ------------------------------------------------------------

In the simple 4-bit shift register, the serial input enters the leftmost
flip-flop and each stored bit moves one position toward the serial output on
the active clock transition.

\begin{whiteboard}

\corresponds{2-8}

\begin{center}
\includegraphics[
    width=0.90\textwidth
]{7_fig_2_8.png}
\end{center}

\[
\boxed{
\text{Serial In}
\rightarrow
A_0
\rightarrow
A_1
\rightarrow
A_2
\rightarrow
A_3
\rightarrow
\text{Serial Out}
}
\]

\end{whiteboard}

% ------------------------------------------------------------
\subsection{Serial and Parallel Transfer}
% ------------------------------------------------------------

Parallel transfer moves multiple bits simultaneously over multiple data lines.

Serial transfer moves one bit at a time over a single data line.

Shift registers can therefore provide conversion between serial and parallel
representations.

\begin{whiteboard}

\[
\boxed{
\text{Parallel transfer: }n\text{ bits at once on }n\text{ lines}
}
\]

\[
\boxed{
\text{Serial transfer: }1\text{ bit at a time on }1\text{ line}
}
\]

Transmitter:

\[
\boxed{
\text{parallel}
\rightarrow
\text{serial}
}
\]

Receiver:

\[
\boxed{
\text{serial}
\rightarrow
\text{parallel}
}
\]

\end{whiteboard}

% ------------------------------------------------------------
\subsection{Bidirectional Shift Register with Parallel Load}
% ------------------------------------------------------------

A more general shift register can:

\begin{itemize}
    \item hold its present contents,
    \item shift right,
    \item shift left,
    \item load data in parallel.
\end{itemize}

Mano's 4-bit design uses one 4-to-1 multiplexer for each D flip-flop.

The selection lines \(S_1\) and \(S_0\) choose which data source reaches each
flip-flop input.

\begin{whiteboard}

\corresponds{2-9}

\begin{center}
\includegraphics[
    height=0.63\textheight,
    keepaspectratio
]{8_fig_2_9.png}
\end{center}

\[
\begin{array}{cc|l}
S_1&S_0&\text{Register Operation}\\
\hline
0&0&\text{No change}\\
0&1&\text{Shift right}\\
1&0&\text{Shift left}\\
1&1&\text{Parallel load}
\end{array}
\]

\end{whiteboard}

The four multiplexer selections have a direct interpretation:

\[
00
\Rightarrow
\text{feedback},
\]

\[
01
\Rightarrow
\text{shift-right source},
\]

\[
10
\Rightarrow
\text{shift-left source},
\]

\[
11
\Rightarrow
\text{parallel input}.
\]

\begin{whiteboard}

\[
\boxed{
00\Rightarrow\text{hold}
}
\]

\[
\boxed{
01\Rightarrow\text{shift right}
}
\]

\[
\boxed{
10\Rightarrow\text{shift left}
}
\]

\[
\boxed{
11\Rightarrow\text{parallel load}
}
\]

\end{whiteboard}

% ============================================================
\section{Binary Counters}
% ============================================================

A counter is a register that passes through a predetermined sequence of states
when input pulses are applied.

A binary counter follows the ordinary binary number sequence.

An \(n\)-bit binary counter contains \(n\) flip-flops and counts from

\[
0
\]

through

\[
2^n-1.
\]

\begin{whiteboard}

\[
\boxed{
\text{Counter}
=
\text{register that follows a state sequence}
}
\]

\[
\boxed{
n\text{-bit binary counter: }
0\rightarrow2^n-1
}
\]

Example:

\[
0000
\rightarrow
0001
\rightarrow
0010
\rightarrow
0011
\rightarrow
0100
\rightarrow
\cdots
\]

\end{whiteboard}

Inspection of the binary counting sequence gives a regular toggle pattern.

The least significant bit changes after every count.

Each higher-order bit changes only when all lower-order bits are 1.

\begin{whiteboard}

\[
\boxed{
A_0\text{ toggles every count}
}
\]

\[
\boxed{
A_i\text{ toggles iff all lower-order bits are }1
}
\]

Example:

\[
0111
\rightarrow
1000
\]

All four bits change because:

\[
A_0=1,
\qquad
A_1A_0=11,
\qquad
A_2A_1A_0=111.
\]

\end{whiteboard}

% ------------------------------------------------------------
\subsection{4-Bit Synchronous Binary Counter}
% ------------------------------------------------------------

JK flip-flops are convenient for binary counters because they toggle when

\[
J=K=1.
\]

In a synchronous counter, every flip-flop receives the same clock.

Combinational logic determines which stages are permitted to toggle on each
clock transition.

\begin{whiteboard}

\corresponds{2-10}

\begin{center}
\includegraphics[
    height=0.62\textheight,
    keepaspectratio
]{9_fig_2_10.png}
\end{center}

\[
\text{Enable}=0
\Rightarrow
\boxed{\text{hold}}
\]

\[
\text{Enable}=1
\Rightarrow
\boxed{\text{count}}
\]

\[
A_0:
\quad
J_0=K_0=\text{Enable}
\]

Higher stages toggle when:

\[
\boxed{
\text{Enable}=1
\text{ and all lower-order bits}=1
}
\]

\end{whiteboard}

The output carry is generated when all stages reach the condition required to
advance an additional counter stage.

It can therefore be used to expand the counter.

% ------------------------------------------------------------
\subsection{Binary Counter with Parallel Load}
% ------------------------------------------------------------

Counters used in digital systems often require more operations than simple
counting.

Mano's counter with parallel load provides:

\begin{itemize}
    \item no change,
    \item increment,
    \item parallel load,
    \item synchronous clear.
\end{itemize}

\begin{whiteboard}

\corresponds{2-11}

\begin{center}
\includegraphics[
    height=0.64\textheight,
    keepaspectratio
]{10_fig_2_11.png}
\end{center}

\[
\begin{array}{ccc|l}
\text{Clear}
&
\text{Load}
&
\text{Increment}
&
\text{Operation}
\\
\hline
0&0&0&\text{No change}\\
0&0&1&\text{Increment by 1}\\
0&1&X&\text{Load }I_0\text{--}I_3\\
1&X&X&\text{Clear to 0}
\end{array}
\]

\end{whiteboard}

The increment operation adds one to the binary number stored in the register.

\begin{whiteboard}

\[
\boxed{
\text{Increment}
=
\text{register contents}+1
}
\]

Priority:

\[
\boxed{
\text{Clear}
>
\text{Load}
>
\text{Increment}
>
\text{Hold}
}
\]

\end{whiteboard}

% ============================================================
\section{Memory Unit}
% ============================================================

A memory unit is a collection of storage cells together with the circuitry
required to transfer information into and out of storage.

Binary information is stored in groups called \emph{words}.

A word is transferred as one unit.

A group of eight bits is called a \emph{byte}.

\begin{whiteboard}

\[
\boxed{
\text{word}
=
\text{group of bits stored/transferred as one unit}
}
\]

\[
\boxed{
1\text{ byte}=8\text{ bits}
}
\]

Examples:

\[
16\text{-bit word}=2\text{ bytes}
\]

\[
32\text{-bit word}=4\text{ bytes}
\]

\end{whiteboard}

% ------------------------------------------------------------
\subsection{Memory Addressing}
% ------------------------------------------------------------

Each word in memory has a unique address.

If a memory contains

\[
2^k
\]

words, then

\[
k
\]

address lines are required to select one word.

If each word contains \(n\) bits, the memory organization can be written as

\[
2^k\times n.
\]

\begin{whiteboard}

\[
\boxed{
2^k\text{ words}
\Rightarrow
k\text{ address lines}
}
\]

\[
\boxed{
\text{memory organization}
=
2^k\text{ words}\times n\text{ bits/word}
}
\]

Address range:

\[
0
\rightarrow
2^k-1
\]

\end{whiteboard}

Memory sizes are commonly expressed using powers of two:

\[
K=2^{10},
\]

\[
M=2^{20},
\]

\[
G=2^{30}.
\]

\begin{whiteboard}

\[
\boxed{
K=2^{10}
}
\qquad
\boxed{
M=2^{20}
}
\qquad
\boxed{
G=2^{30}
}
\]

Examples:

\[
64K=2^{16}
\]

\[
2M=2^{21}
\]

\[
4G=2^{32}
\]

\end{whiteboard}

% ============================================================
\section{Random-Access Memory}
% ============================================================

Random-access memory, or RAM, permits information to be transferred to or
from any selected memory word.

The word is selected by its binary address.

A RAM unit therefore contains:

\begin{itemize}
    \item address lines,
    \item data input lines,
    \item data output lines,
    \item read and write control lines.
\end{itemize}

\begin{whiteboard}

\corresponds{2-12}

\begin{center}
\includegraphics[
    width=0.62\textwidth
]{11_fig_2_12.png}
\end{center}

\[
\boxed{
2^k\text{ words}
\times
n\text{ bits/word}
}
\]

\[
k\text{ address lines}
\]

\[
n\text{ data input lines}
\]

\[
n\text{ data output lines}
\]

\[
\text{Read},
\qquad
\text{Write}
\]

\end{whiteboard}

% ------------------------------------------------------------
\subsection{Write Operation}
% ------------------------------------------------------------

To store a new word in RAM:

\begin{enumerate}
    \item Apply the desired binary address.
    \item Apply the data to the input lines.
    \item Activate the Write control.
\end{enumerate}

\begin{whiteboard}

\[
\boxed{\text{RAM Write}}
\]

\[
1.\ \text{Address}
\]

\[
2.\ \text{Data In}
\]

\[
3.\ \text{Write}
\]

\[
\boxed{
\text{selected word}
\leftarrow
\text{input data}
}
\]

\end{whiteboard}

% ------------------------------------------------------------
\subsection{Read Operation}
% ------------------------------------------------------------

To read a stored word:

\begin{enumerate}
    \item Apply the desired binary address.
    \item Activate the Read control.
\end{enumerate}

The selected memory word is then placed on the output data lines.

Reading does not alter the stored word.

\begin{whiteboard}

\[
\boxed{\text{RAM Read}}
\]

\[
1.\ \text{Address}
\]

\[
2.\ \text{Read}
\]

\[
3.\ \text{Data Out}
\]

\[
\boxed{
\text{Read does not change stored data}
}
\]

\end{whiteboard}

% ============================================================
\section{Read-Only Memory}
% ============================================================

Read-only memory, or ROM, contains permanently established binary information.

A ROM has address inputs and data outputs but does not provide the normal
write capability of RAM.

An \(m\times n\) ROM contains

\[
m
\]

words of

\[
n
\]

bits each.

If

\[
m=2^k,
\]

then \(k\) address lines are needed.

\begin{whiteboard}

\corresponds{2-13}

\begin{center}
\includegraphics[
    width=0.56\textwidth
]{12_fig_2_13.png}
\end{center}

\[
m=2^k
\]

\[
\boxed{
k\text{ address inputs}
\rightarrow
m\text{ stored words}
}
\]

\[
\boxed{
n\text{ output lines}
}
\]

\end{whiteboard}

A ROM output depends only on its current address inputs.

For this reason, Mano classifies ROM as a combinational circuit.

A ROM can implement an input-output relationship specified by a truth table.

\begin{whiteboard}

\[
\boxed{
\text{ROM output}
=
f(\text{address inputs})
}
\]

Therefore:

\[
\boxed{
\text{ROM is a combinational circuit}
}
\]

\[
\boxed{
\text{ROM can implement a truth table}
}
\]

Internal idea:

\[
\boxed{
\text{decoder}
+
\text{OR circuitry}
}
\]

\end{whiteboard}

% ------------------------------------------------------------
\subsection{Types of ROM}
% ------------------------------------------------------------

Mano describes several ways in which ROM information may be established.

\begin{whiteboard}

\[
\boxed{\text{Types of ROM}}
\]

\[
\begin{array}{c|l}
\text{ROM}
&
\text{fixed during manufacture}
\\
\text{PROM}
&
\text{user-programmable once}
\\
\text{EPROM}
&
\text{erasable and reprogrammable}
\\
\text{EEPROM}
&
\text{electrically erasable}
\end{array}
\]

\end{whiteboard}

% ============================================================
\section{Chapter Summary}
% ============================================================

Chapter 2 moves from individual logic gates to reusable digital components.

Decoders and encoders convert between binary codes and one-of-many
representations.

Multiplexers select one data source.

Registers store binary words.

Shift registers move stored information.

Counters advance through ordered binary states.

Memory units provide large collections of individually addressable words.

\begin{whiteboard}

\[
\boxed{
\text{gates}
\rightarrow
\text{digital components}
}
\]

\[
\boxed{
\text{decoder}
\rightarrow
\text{select one output}
}
\]

\[
\boxed{
\text{encoder}
\rightarrow
\text{encode one active input}
}
\]

\[
\boxed{
\text{MUX}
\rightarrow
\text{select one data input}
}
\]

\[
\boxed{
\text{register}
\rightarrow
\text{store bits}
}
\]

\[
\boxed{
\text{shift register}
\rightarrow
\text{move bits}
}
\]

\[
\boxed{
\text{counter}
\rightarrow
\text{sequence through states}
}
\]

\[
\boxed{
\text{memory}
\rightarrow
\text{store addressed words}
}
\]

\end{whiteboard}

\end{document}